% GENERATED FILE. Edit canonical lessons, then run npm run generate:books.
% canonical-source-hash: fnv1a64:c0ee3a93f70af39e
% canonical-lessons: ML-C433-pravesanam, ML-C433-pravesanapariksa, ML-C433-kolsu, ML-C433-semasrrar, ML-C433-silabasu

\chapter{Getting In}
\label{ch:ml-getting-in}
\hlchaptermodality{433}

\begin{chapteropening}
\textbf{By the end of this chapter:} \emph{I can say admission, an entrance exam, a course, a semester and a syllabus.}

Complete the last lesson of chapter 433: I can say admission, an entrance exam, a course, a semester and a syllabus.
\end{chapteropening}

\section[\texorpdfstring{\ml{പ്രവേശനം} (pravēśanaṁ)}{പ്രവേശനം (pravēśanaṁ)}]{\ml{പ്രവേശനം} (pravēśanaṁ) — admission; entry}

\label{lesson:ML-C433-pravesanam}



\begin{quote}
Before the new one: say the Malayalam for a professor, then the Malayalam for a lecture.
\end{quote}

\subsection*{You'll want to know: \ml{പ്രവേശനം}}
\textbf{\ml{പ്രവേശനം}} — \emph{pravēśanaṁ} — \textquotedblleft{}admission; entry\textquotedblright{}.

\begin{quote}
\ml{അവൾക്ക്} \ml{നഗരത്തിലെ} \ml{കോളേജിൽ} \ml{പ്രവേശനം} \ml{കിട്ടി}.
\end{quote}

\begin{quote}
\emph{avaḷkkŭ nagarattile kōḷējil pravēśanaṁ kiṭṭi.}
\end{quote}

\begin{quote}
She got admission to the college in the city.
\end{quote}

On a door: \ml{പ്രവേശനം} \ml{ഇല്ല}, no entry.

\begin{cousinweb}
\textbf{\ml{പ്രവേശനം}} (\emph{pravēśanaṁ}) is borrowed from Sanskrit \emph{praveśana}, \textquotedblleft{}entering, entrance\textquotedblright{} (Monier-Williams).
\end{cousinweb}

\begin{grammarlens}[title={What you've built}]
One more word for giving reasons and opinions.
\end{grammarlens}

\subsection*{Your turn}
\begin{itemize}
\raggedright
  \item \emph{Say it:} \emph{pravēśanaṁ}
  \item \emph{Say it:} \emph{pravēśanaṁ}, once more
  \item \emph{Say it:} \emph{avaḷkkŭ nagarattile kōḷējil pravēśanaṁ kiṭṭi.} — She got admission to the college in the city.
  \item \emph{Recall:} say the Malayalam for a professor, then the Malayalam for a lecture, then say \emph{pravēśanaṁ} again
\end{itemize}

\begin{tcolorbox}[breakable,colback=teal!4,colframe=teal!35!black,title={Before you move on}]
What does \ml{പ്രവേശനം} mean? (\textquotedblleft{}Admission; entry\textquotedblright{}.) Say it once more.
\end{tcolorbox}
\section[\texorpdfstring{\ml{പ്രവേശനപരീക്ഷ} (pravēśanaparīkṣa)}{പ്രവേശനപരീക്ഷ (pravēśanaparīkṣa)}]{\ml{പ്രവേശനപരീക്ഷ} (pravēśanaparīkṣa) — an entrance exam}

\label{lesson:ML-C433-pravesanapariksa}



\begin{quote}
Before the new one: say the Malayalam for a lecture, then the Malayalam for admission.
\end{quote}

\subsection*{You'll want to know: \ml{പ്രവേശനപരീക്ഷ}}
\textbf{\ml{പ്രവേശനപരീക്ഷ}} — \emph{pravēśanaparīkṣa} — \textquotedblleft{}an entrance exam\textquotedblright{}.

\begin{quote}
\ml{അവൻ} \ml{പ്രവേശനപരീക്ഷയ്ക്ക്} \ml{വേണ്ടി} \ml{ഒരു} \ml{വർഷം} \ml{പഠിച്ചു}.
\end{quote}

\begin{quote}
\emph{avan pravēśanaparīkṣaykkŭ vēṇṭi oru varṣaṁ paṭhiccu.}
\end{quote}

\begin{quote}
He studied for a year for the entrance exam.
\end{quote}

\ml{പ്രവേശനം}, entry, and \ml{പരീക്ഷ}, an exam.

\begin{cousinweb}
\textbf{\ml{പ്രവേശനപരീക്ഷ}} (\emph{pravēśanaparīkṣa}) is a Malayalam compound built from two Sanskrit words: \emph{praveśana}, \textquotedblleft{}entering\textquotedblright{}, and \emph{parīkṣā}, \textquotedblleft{}examination, test\textquotedblright{} (Monier-Williams).
\end{cousinweb}

\begin{grammarlens}[title={What you've built}]
One more word for giving reasons and opinions.
\end{grammarlens}

\subsection*{Your turn}
\begin{itemize}
\raggedright
  \item \emph{Say it:} \emph{pravēśanaparīkṣa}
  \item \emph{Say it:} \emph{pravēśanaparīkṣa}, once more
  \item \emph{Say it:} \emph{avan pravēśanaparīkṣaykkŭ vēṇṭi oru varṣaṁ paṭhiccu.} — He studied for a year for the entrance exam.
  \item \emph{Recall:} say the Malayalam for a lecture, then the Malayalam for admission, then say \emph{pravēśanaparīkṣa} again
\end{itemize}

\begin{tcolorbox}[breakable,colback=teal!4,colframe=teal!35!black,title={Before you move on}]
What does \ml{പ്രവേശനപരീക്ഷ} mean? (\textquotedblleft{}An entrance exam\textquotedblright{}.) Say it once more.
\end{tcolorbox}
\section[\texorpdfstring{\ml{കോഴ്സ്} (kōḻsŭ)}{കോഴ്സ് (kōḻsŭ)}]{\ml{കോഴ്സ്} (kōḻsŭ) — a course of study}

\label{lesson:ML-C433-kolsu}



\begin{quote}
Before the new one: say the Malayalam for admission, then the Malayalam for an entrance exam.
\end{quote}

\subsection*{You'll want to know: \ml{കോഴ്സ്}}
\textbf{\ml{കോഴ്സ്}} — \emph{kōḻsŭ} — \textquotedblleft{}a course of study\textquotedblright{}.

\begin{quote}
\ml{അവൾ} \ml{ഒരു} \ml{കമ്പ്യൂട്ടർ} \ml{കോഴ്സ്} \ml{ചെയ്യുന്നു}.
\end{quote}

\begin{quote}
\emph{avaḷ oru kampyūṭṭar kōḻsŭ ceyyunnu.}
\end{quote}

\begin{quote}
She is doing a computer course.
\end{quote}

One does a course: \ml{കോഴ്സ്} \ml{ചെയ്യുക}.

\begin{cousinweb}
\textbf{\ml{കോഴ്സ്}} (\emph{kōḻsŭ}) was borrowed from English \emph{course}.
\end{cousinweb}

\begin{grammarlens}[title={What you've built}]
One more word for giving reasons and opinions.
\end{grammarlens}

\subsection*{Your turn}
\begin{itemize}
\raggedright
  \item \emph{Say it:} \emph{kōḻsŭ}
  \item \emph{Say it:} \emph{kōḻsŭ}, once more
  \item \emph{Say it:} \emph{avaḷ oru kampyūṭṭar kōḻsŭ ceyyunnu.} — She is doing a computer course.
  \item \emph{Recall:} say the Malayalam for admission, then the Malayalam for an entrance exam, then say \emph{kōḻsŭ} again
\end{itemize}

\begin{tcolorbox}[breakable,colback=teal!4,colframe=teal!35!black,title={Before you move on}]
What does \ml{കോഴ്സ്} mean? (\textquotedblleft{}A course of study\textquotedblright{}.) Say it once more.
\end{tcolorbox}
\section[\texorpdfstring{\ml{സെമസ്റ്റർ} (semasṟṟar)}{സെമസ്റ്റർ (semasṟṟar)}]{\ml{സെമസ്റ്റർ} (semasṟṟar) — a semester}

\label{lesson:ML-C433-semasrrar}



\begin{quote}
Before the new one: say the Malayalam for an entrance exam, then the Malayalam for a course of study.
\end{quote}

\subsection*{You'll want to know: \ml{സെമസ്റ്റർ}}
\textbf{\ml{സെമസ്റ്റർ}} — \emph{semasṟṟar} — \textquotedblleft{}a semester\textquotedblright{}.

\begin{quote}
\ml{അടുത്ത} \ml{സെമസ്റ്റർ} \ml{ജൂണിൽ} \ml{തുടങ്ങും}.
\end{quote}

\begin{quote}
\emph{aṭutta semasṟṟar jūṇil tuṭaṅṅuṁ.}
\end{quote}

\begin{quote}
The next semester will begin in June.
\end{quote}

Plural: \ml{സെമസ്റ്ററുകൾ}.

\begin{cousinweb}
\textbf{\ml{സെമസ്റ്റർ}} (\emph{semasṟṟar}) was borrowed from English \emph{semester}.
\end{cousinweb}

\begin{grammarlens}[title={What you've built}]
One more word for giving reasons and opinions.
\end{grammarlens}

\subsection*{Your turn}
\begin{itemize}
\raggedright
  \item \emph{Say it:} \emph{semasṟṟar}
  \item \emph{Say it:} \emph{semasṟṟar}, once more
  \item \emph{Say it:} \emph{aṭutta semasṟṟar jūṇil tuṭaṅṅuṁ.} — The next semester will begin in June.
  \item \emph{Recall:} say the Malayalam for an entrance exam, then the Malayalam for a course of study, then say \emph{semasṟṟar} again
\end{itemize}

\begin{tcolorbox}[breakable,colback=teal!4,colframe=teal!35!black,title={Before you move on}]
What does \ml{സെമസ്റ്റർ} mean? (\textquotedblleft{}A semester\textquotedblright{}.) Say it once more.
\end{tcolorbox}
\section[\texorpdfstring{\ml{സിലബസ്} (silabasŭ)}{സിലബസ് (silabasŭ)}]{\ml{സിലബസ്} (silabasŭ) — a syllabus}

\label{lesson:ML-C433-silabasu}



\begin{quote}
Before the new one: say the Malayalam for a course of study, then the Malayalam for a semester.
\end{quote}

\subsection*{You'll want to know: \ml{സിലബസ്}}
\textbf{\ml{സിലബസ്}} — \emph{silabasŭ} — \textquotedblleft{}a syllabus\textquotedblright{}.

\begin{quote}
\ml{ഇത്തവണ} \ml{സിലബസ്} \ml{വളരെ} \ml{വലുതാണ്}.
\end{quote}

\begin{quote}
\emph{ittavaṇa silabasŭ vaḷare valutāṇŭ.}
\end{quote}

\begin{quote}
This time the syllabus is very long.
\end{quote}

Plural: \ml{സിലബസുകൾ}.

\begin{cousinweb}
\textbf{\ml{സിലബസ്}} (\emph{silabasŭ}) was borrowed from English \emph{syllabus}.
\end{cousinweb}

\begin{grammarlens}[title={What you've built}]
One more word for giving reasons and opinions.
\end{grammarlens}

\subsection*{Your turn}
\begin{itemize}
\raggedright
  \item \emph{Say it:} \emph{silabasŭ}
  \item \emph{Say it:} \emph{silabasŭ}, once more
  \item \emph{Say it:} \emph{ittavaṇa silabasŭ vaḷare valutāṇŭ.} — This time the syllabus is very long.
  \item \emph{Recall:} say the Malayalam for a course of study, then the Malayalam for a semester, then say \emph{silabasŭ} again
\end{itemize}

\begin{tcolorbox}[breakable,colback=teal!4,colframe=teal!35!black,title={Before you move on}]
What does \ml{സിലബസ്} mean? (\textquotedblleft{}A syllabus\textquotedblright{}.) Say it once more.
\end{tcolorbox}
